\chapter{Conclusions}
\label{cap:conclusions}
\let\cleardoublepage\clearpage

\section{Conclusions}
\label{sec:conclusions}

This report has presented the design, implementation, and validation of a comprehensive platform for remote laboratory access, addressing the growing need for flexible and accessible laboratory resources in educational and research environments. The solution successfully bridges the gap between traditional hands-on experimentation and the requirements of modern remote learning and collaboration.

The developed platform successfully fulfills all primary objectives established at the project's inception. The implementation of a robust web \ac{api} provides comprehensive management capabilities for users, laboratories, and hardware resources. The intuitive web application interface ensures accessibility for users with varying technical backgrounds, while the secure authentication and authorization mechanisms, coupled with role-based access control, guarantee system integrity and appropriate resource allocation.

The well-designed database system ensures reliable data persistence, supporting the platform's scalability requirements. Furthermore, the implementation of remote manipulation capabilities of laboratory devices enables real-time interaction with experimental equipment, maintaining the essential hands-on experience that characterizes quality laboratory work.

Beyond its technical accomplishments, this platform marks a significant advancement in the transformation of experimental education. By enabling remote access to laboratory resources, the project democratizes hands-on learning, allowing students and researchers from diverse regions and socioeconomic backgrounds to participate in high-quality laboratory activities without geographical constraints.

The flexibility and scalability of the system open new perspectives for pedagogical innovation, supporting hybrid and collaborative teaching methodologies. It's potential for integration with other educational platforms and for fostering international partnerships further amplifies the project's impact, establishing it as a reference for future initiatives in remote education and research. The system's sustainability and capacity for evolution ensure that it will continue to address emerging challenges in experimental learning.

The development of this project provided an invaluable opportunity to apply knowledge acquired across multiple domains of the academic curriculum, including software engineering, web development, database management, networking, and security. Throughout the process, new skills and insights were also gained, particularly in the areas of system integration, remote device management, and collaborative software design. This interdisciplinary experience not only reinforced existing competencies but also fostered significant personal and academic growth.

\section{Future Work}
\label{sec:future_work}

While the current implementation successfully addresses the primary objectives, several areas present opportunities for future enhancement. The expansion of the testing framework to include comprehensive load testing will provide valuable insights into system performance under high-usage conditions. This analysis will inform optimization strategies and guide infrastructure scaling decisions.

Future development efforts should focus on expanding hardware compatibility to support a broader range of laboratory equipment types. Additionally, the implementation of advanced monitoring and analytics capabilities could provide valuable usage insights and facilitate resource optimization. The integration of collaborative features, such as shared experimentation sessions and real-time collaboration tools, would further enhance the platform's educational value.